% !PRINTED 366
\weberpage{366}
Nun haben wir die beiden Complexe:

\weq{4}{x_2 z_2,\ x_3 z_3,\ \dots}
\weq{4\text{b}}{x_2 x_3,\ z_2 z_3,\ \dots{},}

und da $x_2 x_3 y_2 y_3$ ein syzygetisches Quadrupel sind, so enthält die Gruppe (4 b) auch das Paar $y_2 y_3$, und folglich sind auch $y_2 y_3 z_2 z_3$ ein syzygetisches Quadrupel. Daraus folgt weiter, dass $y_2 z_2$ und $y_3 z_3$ in denselben Complex gehören. Da man dieselbe Betrachtung wie für $x_3 y_3$ auch für die Paare $x_4 y_4, x_5 y_5, x_6 y_6$ durchführen kann, so lassen sich hiernach die Complexe (4), (5) vollständig bilden, und sie erhalten den Ausdruck:

\weq{4}{x_1 z_1,\ x_2 z_2,\ x_3 z_3,\ x_4 z_4,\ x_5 z_5,\ x_6 z_6}
\weq{5}{y_1 z_1,\ y_2 z_2,\ y_3 z_3,\ y_4 z_4,\ y_5 z_5,\ y_6 z_6.}

Hierin bilden $z_1 z_2$ ein Paar des Complexes (2) und $z_3, z_4, z_5, z_6$, die ein azygetisches System bilden, kommen alle in dem Complex (3) vor, in dem keine zwei gepaart sind.

Da wir, wie vorhin gezeigt, nach dem Typus (2), (3), (4), (5) aus dem willkürlich angenommenen Complex (1) alle überhaupt existirenden Complexe ableiten können, so ergiebt sich der Satz:

6. \gesperrt{Irgend zwei Complexe haben entweder ein syzygetisches Quadrupel oder ein azygetisches System von sechs Doppeltangenten gemein.}

Zwei Complexe, die ein syzygetisches Quadrupel gemein haben, wollen wir ein \gesperrt{syzygetisches Complexpaar} nennen. Zu jedem solchen Paare giebt es einen dritten Complex, der dasselbe Quadrupel enthält. Drei solche Complexe nennen wir ein \gesperrt{syzygetisches Complextripel} [z.\,B. (1), (2), (3)].

Ebenso nennen wir zwei Complexe der zweiten Art, d.\,h. solche, die sechs azygetische Elemente gemein haben, ein \gesperrt{azygetisches Complexpaar}.

Jedes azygetische Complexpaar wird gleichfalls durch einen bestimmten Complex, der mit jedem der beiden Complexe ein azygetisches Paar bildet, zu einem Tripel ergänzt [wie (1), (4), (5)]. Ein solches nennen wir ein \gesperrt{azygetisches Complextripel}.

In einem syzygetischen Tripel kommen alle 28 Doppeltangenten vor, in einem azygetischen nur 18.
